\subsection*{集団力学と心理学}

\addcontentsline{toc}{subsection}{2 集団力学と心理学}

ソフトウェア工学の仕事は、多くの場合チームで行われる。チームには、開発者、テスト担当、運用担当、顧客、利用者、管理者、規制当局、外部委託先などが関わる。したがって、技術者には、他者と協力し、対立を建設的に扱い、情報を共有し、継続的に改善する力が求められる。

\subsubsection*{チーム・集団で働く力学}

\addcontentsline{toc}{subsubsection}{2.1 チーム・集団で働く力学}

成果を出すチームは、凝集性を持ち、協力的で、正直で、焦点の合った雰囲気を持つ。個人の目標とチームの目標がそろっていると、メンバーは共有成果に自然に責任と所有感を持つ。

よいチームでは、知的に正直であること、知らないことを認めること、誤りを認めること、責任・報酬・作業負荷を公平に分けることが重視される。文書、ソースコード、会話で明確に伝えることで、情報が全員に届く。レビューでは人ではなく成果物に焦点を当て、非個人的かつ建設的に行う。

ソフトウェアは多様な領域に入り込むため、技術者は多分野の環境で働けなければならない。スマートフォン、自動車、医療機器、業務システムなど、適用領域が変われば、関係者の価値観やリスクも変わる。

\subsubsection*{個人の認知}

\addcontentsline{toc}{subsubsection}{2.2 個人の認知}

個人の認知とは、人が情報を理解し、記憶し、問題を分解し、解決策を考える仕組みである。ソフトウェアは抽象度が高いため、個人の認知限界は問題解決に大きく影響する。

\begin{center}
\small
\par\smallskip\noindent\refstepcounter{table}\label{tab:ch14-04}表 \thetable: 個人の認知の整理\par\smallskip
表 \thetable は、個人の認知の整理を比較・確認しやすい形で整理したものである。\par\smallskip
\begin{longtable}{p{0.27\linewidth}p{0.68\linewidth}}
\toprule
認知を妨げる要因 & 説明 \\
\midrule
\endfirsthead
\toprule
認知を妨げる要因 & 説明 \\
\midrule
\endhead
知識不足 & 対象領域、技術、標準、制約を知らないために正しく判断できない。 \\
\midrule
無意識の前提 & 自分では当然と思っている条件が、実際には成り立たない。 \\
\midrule
データ量の多さ & 情報が多すぎて、重要な点を見落とす。 \\
\midrule
失敗への恐れ & 評価や責任を恐れ、相談や試行を避ける。 \\
\midrule
文化 & 組織文化や適用領域の文化が、質問や発言を妨げる。 \\
\midrule
表現力不足 & 問題を言語化・図式化できず、共有できない。 \\
\midrule
雰囲気 & 発言しにくい職場環境が問題発見を遅らせる。 \\
\midrule
感情状態 & 疲労、不安、怒りなどが判断に影響する。 \\
\midrule
\bottomrule
\end{longtable}
\normalsize
\end{center}

\par\smallskip
\begin{center}
\centering
\IfFileExists{assets/generated_figures/ch14/fig-ch14-07.png}{\includegraphics[width=0.96\linewidth]{assets/generated_figures/ch14/fig-ch14-07.png}}{\fbox{\parbox[c][48mm][c]{0.9\linewidth}{\centering 画像生成待ち\\個人の認知限界と対策を示す図}}}
\par\smallskip
\refstepcounter{figure}\label{fig:ch14-07}図 \thefigure: 個人の認知限界と対策を示す図
\end{center}
図 \thefigure は、個人の認知限界と対策を示す図を視覚的に整理したものである。
\par\smallskip

\subsubsection*{問題の複雑さへの対応}

\addcontentsline{toc}{subsubsection}{2.3 問題の複雑さへの対応}

多くのソフトウェア工学問題は、全体を一度に扱うには複雑すぎる。そのため、チーム作業と問題分解が必要になる。問題分解とは、大きな問題を小さな部分に分け、それぞれを理解・解決し、最後に統合する考え方である。

チームでは、メンバーごとの知識、経験、視点、創造性が相補的に働く。ペアプログラミングやコードレビューは、複雑な問題に対するチームでの対処例である。

\subsubsection*{利害関係者との相互作用}

\addcontentsline{toc}{subsubsection}{2.4 利害関係者との相互作用}

ソフトウェア工学の成功は、利害関係者との良好な相互作用に依存する。利害関係者は、ソフトウェアライフサイクルのすべての段階で、支援、情報、フィードバックを提供する。初期段階では、すべての利害関係者を特定し、製品がどのように影響するかを理解することが重要である。

アジャイル開発では、利害関係者の関与が特に重要である。開発中、利害関係者は仕様や早期版へのフィードバック、優先度変更、詳細要求の明確化を提供する。保守中も、進化する要求や問題報告を通じて、ソフトウェアの拡張と改善に関わる。

\par\smallskip
\begin{center}
\centering
\IfFileExists{assets/generated_figures/ch14/fig-ch14-08.png}{\includegraphics[width=0.96\linewidth]{assets/generated_figures/ch14/fig-ch14-08.png}}{\fbox{\parbox[c][48mm][c]{0.9\linewidth}{\centering 画像生成待ち\\利害関係者関与のライフサイクル図}}}
\par\smallskip
\refstepcounter{figure}\label{fig:ch14-08}図 \thefigure: 利害関係者関与のライフサイクル図
\end{center}
図 \thefigure は、利害関係者関与のライフサイクル図を視覚的に整理したものである。
\par\smallskip

\subsubsection*{不確実性と曖昧さへの対応}

\addcontentsline{toc}{subsubsection}{2.5 不確実性と曖昧さへの対応}

ソフトウェア技術者は、不確実性と曖昧さを扱わなければならない。不確実性とは、必要な情報が不足している状態である。曖昧さとは、同じ言葉や状況が複数の意味に解釈できる状態である。

不確実性が知識不足に由来する場合は、教科書、専門誌、標準、論文、公式資料を読む、利害関係者に面談する、チームメイトや専門家に相談することで解決できることがある。容易に解消できない不確実性は、プロジェクトリスクとして扱い、見積りや価格、計画に反映する。

\par\smallskip
\begin{center}
\centering
\IfFileExists{assets/generated_figures/ch14/fig-ch14-09.png}{\includegraphics[width=0.96\linewidth]{assets/generated_figures/ch14/fig-ch14-09.png}}{\fbox{\parbox[c][48mm][c]{0.9\linewidth}{\centering 画像生成待ち\\不確実性への対応フロー}}}
\par\smallskip
\refstepcounter{figure}\label{fig:ch14-09}図 \thefigure: 不確実性への対応フロー
\end{center}
図 \thefigure は、不確実性への対応フローを視覚的に整理したものである。
\par\smallskip

\subsubsection*{公平性・多様性・包摂性への対応}

\addcontentsline{toc}{subsubsection}{2.6 公平性・多様性・包摂性への対応}

公平性、多様性、包摂性は、チームの力学に影響する。地理的に分散したチームやコミュニケーション頻度が低いチームでは、一回一回の接触の重要性が高まる。時差があると口頭でのやり取りが減り、文化的な誤解も起こりやすい。

ソフトウェア開発では国際的な外部委託や分散開発が多いため、多文化環境は特に一般的である。プロジェクトを成功させるには、文化的・社会的規範の違いへの寛容さ、リーダーシップによる支援、より頻繁なコミュニケーション、可能な範囲での対面または同期的な会議が有効である。

性別による偏りもソフトウェア産業の課題である。より広い採用戦略、具体的で測定可能な評価基準、報酬決定手続きの透明化は、不平等を減らすための手段になり得る。

\par\smallskip
\begin{center}
\centering
\IfFileExists{assets/generated_figures/ch14/fig-ch14-10.png}{\includegraphics[width=0.96\linewidth]{assets/generated_figures/ch14/fig-ch14-10.png}}{\fbox{\parbox[c][48mm][c]{0.9\linewidth}{\centering 画像生成待ち\\チーム人数とコミュニケーション経路の増加図}}}
\par\smallskip
\refstepcounter{figure}\label{fig:ch14-10}図 \thefigure: チーム人数とコミュニケーション経路の増加図
\end{center}
図 \thefigure は、チーム人数とコミュニケーション経路の増加図を視覚的に整理したものである。
\par\smallskip

\subsubsection*{プレゼンテーション技能}

\addcontentsline{toc}{subsubsection}{3.4 プレゼンテーション技能}

ソフトウェア技術者は、ライフサイクルの多くの場面でプレゼンテーション技能を使う。要求レビューで顧客に要求を説明する、設計レビューで構造を説明する、製品ウォークスルーを行う、利用者や運用担当へ教育する、といった場面である。

よい発表は、技術情報を一方的に話すだけではない。相手の前提知識に合わせて概念を整理し、意思決定に必要な情報を示し、質問やフィードバックを引き出す。発表で使った知識は、スライド、技術メモ、白書、ナレッジ記事などとして保存し、組織の知識資産にすることが望ましい。

\par\smallskip
\begin{center}
\centering
\IfFileExists{assets/generated_figures/ch14/fig-ch14-11.png}{\includegraphics[width=0.96\linewidth]{assets/generated_figures/ch14/fig-ch14-11.png}}{\fbox{\parbox[c][48mm][c]{0.9\linewidth}{\centering 画像生成待ち\\コミュニケーション技能の関係図}}}
\par\smallskip
\refstepcounter{figure}\label{fig:ch14-11}図 \thefigure: コミュニケーション技能の関係図
\end{center}
図 \thefigure は、コミュニケーション技能の関係図を視覚的に整理したものである。
\par\smallskip
